\section{Transport from the modular block to the incidence quotient}
\label{sec:incidence-transport-new}
The central APP block and face--vertex incidence provide
two realizations of the same involution. Their relation is expressed through
a rectangular intertwiner, before any metric interpretation
is made. The adjacent diagonal congruence
\[
 k^2\equiv(k+1)^2\pmod9
 \quad\Longleftrightarrow\quad 2k+1\equiv0\pmod9
\]
selects \(k=4\) in the declared window. The corresponding block is
\[
 B_c=\begin{pmatrix}7&2\\2&7\end{pmatrix}
 =7I_2+2R_c=9P_++5P_-,\qquad
 R_c=\begin{pmatrix}0&1\\1&0\end{pmatrix},\quad
 P_\pm=\frac{I_2\pm R_c}{2}.
\]
The difference between the two eigenvalues is four. Incidence
transport then determines how this difference is preserved
when sector multiplicities change.

The cube faces are indexed by \((i,\epsilon)\), with
\(i\in\{1,2,3\}\) and \(\epsilon\in\{-1,1\}\); its vertices are
indexed by \(\nu\in\{-1,1\}^3\). Define
\[
 M_{(i,\epsilon),\nu}=\mathbf1_{\{\nu_i=\epsilon\}}.
\]
Let \(R_F\) be face opposition and \(R_V\) vertex inversion
\(\nu\mapsto-\nu\). Entrywise, one verifies
\[
 \mathbf1_{\{\nu_i=-\epsilon\}}
 =\mathbf1_{\{(-\nu)_i=\epsilon\}},\qquad R_FM=MR_V.
\]
The matrix \(M\) therefore transports the full involution.

\begin{theorem}[Incidence descent of the central block]
Let \(B_F=7I_6+2R_F\) and \(B_V=7I_8+2R_V\). Then
\[
 B_FM=MB_V.
\]
If \(u=(1,1,1,1,1,1)^{\mathsf T}\),
\(P_u=uu^{\mathsf T}/6\) and \(P_-=(I_6-R_F)/2\), the quotient
\[
 Q_\square=\RR^6/\ker(MM^{\mathsf T})
 \simeq\RR u\oplus E_-
\]
has dimension \(1+3\), and the induced operator is
\[
 \overline B_F=9P_u+5P_-.
\]
\end{theorem}
\begin{proof}
The first identity follows by multiplying \(R_FM=MR_V\) by two
and adding \(7M\). Each face contains four vertices; two adjacent faces
share two, and opposite faces share none. This count gives,
entrywise,
\[
 MM^{\mathsf T}=12P_u+4P_-.
\]
The first projector has rank one and the second rank three; their images
are orthogonal. The eigenvalues are \(12,4,0\), with multiplicities
\(1,3,2\). Opposition acts as \(+1\) on \(\RR u\) and as
\(-1\) on \(E_-\), and preserves the kernel of the Gram matrix. Thus,
\(B_F\) descends to the quotient and acts by nine and five on those two
summands. Its characteristic polynomial is
\((\lambda-9)(\lambda-5)^3\).
\end{proof}

The temporal assignment to the axis \(\RR u\) fixes the signed form
\(J_L=P_--P_u\). On the quotient, \(I_Q=P_u+P_-\), so
\[
 J_L=\frac{7I_Q-\overline B_F}{2}.
\]
This equality relates the form of signature \((3,1)\), with the
stated temporal convention, to the transported modular operator. The
initial Gram matrix retains its positivity; the signature is introduced
in the subsequent realization through the sign assigned to each sector.
The dilations nine and five remain those of the block: the affine
identity determines the form and does not turn the block into an isometry.

\subsection{Arithmetic aggregation, lattice norm and route functional}
Preservation of the integer before reduction becomes visible by summing
the two sheets. The modular aggregation matrices are
\[
 M_\Pi=\begin{pmatrix}4&5&6\\5&4&6\\6&6&9\end{pmatrix},\qquad
 M_\Sigma=\begin{pmatrix}5&6&4\\6&4&5\\4&5&6\end{pmatrix}.
\]
Their sums are fifty-one and forty-five. Upon reconstructing the
nine elements of each block, the integer and its carry satisfy
\[
 2025-810=(459-405)+9(174-45)=54+9\cdot129.
\]
The residue aggregate and the quotient constitute two coordinates of this
equality. In the subsequent lattice construction, the value fifty-four
appears as the norm of the marked chamber
\(v_\alpha=4\rho_1+\rho_2+\cdots+\rho_{12}\), for orthogonal
vectors \(\rho_j^2=2\). In the material realization it appears as
the evaluation of the route functional \(D_1\) of the profile \((80,54,6)\).
Each occurrence retains its originating map: aggregation with carry,
lattice form and oriented route evaluation. The equality of their
values acquires structural content through these maps, whose
domains and proofs are developed in their respective sections.
